\documentclass{article}
\usepackage[T2A]{fontenc}
\usepackage[utf8]{inputenc}   
\usepackage[english, russian]{babel}

% Set page size and margins
% Replace `letterpaper' with`a4paper' for UK/EU standard size
\usepackage[a4paper,top=2cm,bottom=2cm,left=2cm,right=2cm,marginparwidth=1.75cm]{geometry}

\usepackage{amsmath}
\usepackage{graphicx}
\usepackage[colorlinks=true, allcolors=blue]{hyperref}
\usepackage{amsfonts}
\usepackage{amssymb}
% \usepackage[left=1cm,right=1cm,top=1cm,bottom=1cm]{geometry}
\usepackage{hyperref}
\usepackage{seqsplit}
\usepackage[dvipsnames]{xcolor}
\usepackage{enumitem}
\usepackage{algorithm}
\usepackage{algpseudocode}
\usepackage{algorithmicx}
\usepackage{mathalfa}
\usepackage{mathrsfs}
\usepackage{dsfont}
\usepackage{caption,subcaption}
\usepackage{wrapfig}
\usepackage[stable]{footmisc}
\usepackage{indentfirst}
\usepackage{rotating}
\usepackage{pdflscape}
\usepackage{listings}

\usepackage{minted}

\usepackage{MnSymbol,wasysym}

\title{Домашнее задание 5}
\author{Лиза Фоменко}
\date{}

\begin{document}

\maketitle



\section*{\textbf{Задание 1}} \textbf{\textit{Оцените сложность алгоритма Divide, приведенного на странице 19 книги Дасгупты.}}

Сначала выпишем алгоритм

\begin{verbatim}
0 def Divide(x: int (n bit), y: int >= 1 ) -> q, r:
1     if x==0: 
2         return 0,0
3 
4     q, r = Divide(⌊x/2⌋, y)
5     q = 2q
6     r = 2r
7     if x[-1] == 1: 
8         r = r + 1
9     if r >= y: 
10         r = r - y
11         q = q + 1
12    return q, r
\end{verbatim}

Рассмотрим алгоритм построчно. за $\Theta(1)$ мы проверяем, равен ли x 0. 
Затем мы рекурсивно вызываем функцию $Devide$ для $\left \lfloor{\frac{x}{2}}\right \rfloor $ и $y$. Заметим, что (как оговаривалось на лекции), деление числа на 2 реализовано за $O(1)$, так как мы сдвигаем число на один разряд вправо. Каждый новый вызов $Devide$ мы сдвигаем $x$ на еще один разряд вправо. Так как по условию $x$ - n-битное число, то таких сдвигов будет ровно n. Значит, функция $Devide$ будет рекурсивно вызвана $n$ раз.

Далее оценим тело самой функции.

В строках 4 и 5 происходит умножение на 2 -- то есть левый сдвиг за $O(1)$ каждое. Итого $O(1)$ за две строки.

В строке 7 мы проверяем число на нечетность: для этого нужно последний бит числа $x$, $x[-1]$, сравнить с 1. Это атомарная операция, выполняется за $O(1)$. При необходимости мы выполняем операцию сложения $r = r+1$ за $O(n)$ в строке 8(мне кажется все-таки за n, а не 1, потому что мы принимаем за $O(1)$ атомарные операции, а число $r$ имеет не более $n$ разрядов).

В строках 10-12 также выполняется одна проверка n-битных чисел (за $O(n)$ в атомарных операциях), а далее следуют две операции сложения и вычитания за $O(n)$ (так как $y$ n-битовое число) 

Тогда все тело функции $Devide$ имеет асимптотику $O(n)$. Так как $Devide$ вызывается $n$ раз, то общее время работы алгоритма $O(n^2)$.



\bigskip

\section*{\textbf{Задание 2}} \textbf{\textit{Предложите алгоритм возведения $n$-битовых чисел в $n$-битовую степень по $n$-битовому модулю, оцените его сложность.
}}



\bigskip


\section*{\textbf{Задание 3}} \textbf{\textit{Кто-то очень умный пришел и сказал, что умеет возводить числа из $n$ бит в квадрат за линейное время. Докажите, что тогда можно перемножать и разные числа за $O$ от максимума их длин.}}

Буду пользоваться тем, что $(a + b)^2 = a^2 + 2ab + b^2$.

Выразим отсюда $ab = \frac{(a + b)^2 - a^2 - b^2}{2}$

Теперь исследуем каждую операцию
\begin{itemize}
    \item сложение $a+b$ - сложение двух n-битных чисел, работает за $O(n)$, где n - макимсальная из длин чисел
    \item возведение в квадрат $(a+b)^2$ работает за $O(n)$ согласно кому-то оч умному
    \item возведение n-битных чисел $a, b$ в квадрат $a^2, b^2$ тоже работает за $O(n)$, n - длина числа
    \item осталось два вычитания n-битных чисел за $O(n)$, где n - макиммальная длина числе
    \item деление на 2 - это сдвиг вправо на разряд, за $O(1)$
\end{itemize}

Итого $ab$ можно найти за $O(n)$, где n - максимальная из их длин

\bigskip


\section*{\textbf{Задание 4}} \textbf{\textit{Докажите, что НОД(a, b) = НОД(a - b, b)}}

Пусть мы знаем, что НОД(a, b) = k
Тогда a = kr, b = km

Распишем a - b = kr - km = k(r-m)
Тогда мы ищем НОД(k(r-m), km)

Пусть НОД(k(r-m), km) = l, тогда по условию l>k (мы видим, что k является делителем k(r-m) и km, но по предположению не является их НОДом, тогда k<l). Это значит, что мы можем записать два выражения: 

k(r-m) = ol = a - b;  km = il = b

Из этого равенства мы можем найти a: a = a - b + b = ol + il = l(o+i). Тогда мы получили, что a = l(o+i) и b = il, то есть l является делителем a и делителем b, при этом l>k по предположению, значит, k не может быть НОД(a, b), что протеворечит изначальному условию.

Тогда l = k и НОД(a, b) = НОД(a-b, b) = k


\bigskip


\section*{\textbf{Задание 5}} \textbf{\textit{На доске написан набор положительных целых чисел. За один ход можно взять любые два числа и вычесть из большего меньшее. Процесс останавливается, когда все числа становятся одинаковыми. Докажите, что этот процесс всегда остановится. Какие числа останутся в результате?}}

Докажем, что процесс остановится. 
При случайном выборе $X_1, X_2$, пусть без ограничения общности $X_1 > X_2$. За наш шаг мы меняем $X_1 = X_1 - X_2$. Подозрительно похоже на предыдущий пункт, будем держать в уме, что НОД($X_1, X_2$) от этой операции не меняется. Кстати интересно, что мы выбирали $X_1, X_2$ без ограничения общности, тогда за все попарные НОД между числами не изменятся в ходе наших операций. 

Докажем детерменированность: $X_1 = X_1 - X_2$, при этом $X_1$ остается положительным числом после манипуляции. Так как мы из большего положительного числа всегда вычитаем меньшее положительное число, то мы всегда будем получать положительные числа и никогда не дойдем до отрицательных, поэтому процесс остановится. Остановится он тогда, когда все числа будут одинаковыми и мы не сможем выбрать меньшее из них.

Тогда нужно понять, какое именно число (ну по факту это число будет повторяться n раз, n - размер исходной последовательности) останется. Тут можно вопспользоваться тем, что НОД каждой пары чисел не изменяется после проивзеденной операции: тогда и НОД всего множества чисел не изменится. Тогда, оперируя двумя фактами, все числа равны и НОД множества не изменился с изначального, выводим, что в результате каждое число будет заменено на НОД($X_1, X_2, ... X_n$)

\bigskip



\section*{\textbf{Задание 6}} \textbf{\textit{Оцените асимптотику рекурренты $T(n) = T(\alpha n) + T((1 - \alpha) n)) + \Theta(n)$ для $\alpha \in (0, 1)$.}}

Эту задачу будем решать с помощью дерева рекурции. Пусть (для удобства), без ограничения общности, $\alpha < 1 - \alpha$. 

Докажем с помощью дерева рекурсии. Наше дерево бинарное, но не обязательно полное, его высота: от $\log_{\frac{1}{\alpha}} n$ (самый близкий к корню, мы его получим по ветке $\alpha n$ -> $\alpha n$ -> $\alpha n$ ...) до $\log_{\frac{1}{1-\alpha}} n$ (самый далекий, его получим по ветке $(1-\alpha) n$ -> $(1-\alpha) n$ -> $(1-\alpha) n$ ...). Логика моего решения будет такая:

\begin{verbatim}
                   n
               /       \                -> n
           an          (1-a)n           -> an + (1-a)n = n
        /        \   /       \          -> 
     a^2n   a(1-a)n a(1-a)n  (1-a)^2n    -> a^2n + 2a(1-a)n +(1-a)^2n = n
.......................................
 / \  / \ / \ / \ / \ / \ / \ / \ / \     # level log{1/a} n
 O(1)  
.......................................
                                   / \     # level log{1/(1-a)} n
                                O(1) O(1)
\end{verbatim}

Итак, а каждом уровне до $\log_{\frac{1}{\alpha}} n$ проводится ровно n операций, которые входят в f(n). На уровнях с $\log_{\frac{1}{\alpha}} n$ до $\log_{\frac{1}{1-\alpha}} n$ листья выполняются за $O(1)$, а шагов из f(n) выполняется не более ($\leq n$), так как рекурсия из листьев не вызывается. Тогда мы можем оценить число шагов алгоритма как: $n \log_{\frac{1}{\alpha}} n  \leq T(n) \leq n \log_{\frac{1}{1-\alpha}} n$, а, как мы доказывали раньше (типа в 1 дз что ли), это означает, что $T(n) = \Theta (n \log n)$ 

(Это просто обобщение задачи из предыдущей домашки, 2пункт5, которая еще разбиралась на семинаре)

\bigskip



\section*{\textbf{Задание 7}} \textbf{\textit{Оцените трудоемкость алгоритма, разбивающего задачу размера $n$ на $n$ подзадач размера $\frac{n}{2}$, используя для этого $\Theta(n)$ операций.}}

Попробуем записать так: $T(n) = n T(\frac{n}{2}) + \Theta(n)$

Нарисуем дерево рекурсии. Оно будет полным.
\begin{verbatim}
                   n                       -> n
         /    /    |    \    \             -> 
       n/2  n/2   n/2   n/2   n/2          -> n*n/2 = n^2/2
   /  |  \   |  \    \   |  \   \   \  \    -> 
 n/4 n/4 n/4 n/4 n/4 n/4 n/4 n/4 n/4 n/4 n/4  -> n * n/2 * n/4 = n^3 / 2^3
.......................................
 / \  / \ / \ / \ / \ / \ / \ / \ / \     # 
 O(1)  O(1) O(1) O(1) O(1) O(1) O(1) O(1) O(1) 

\end{verbatim}

Попробуем по порядку найти зависимость.

На первом уровне мы делим задачу размера n на n задач размера n/2, тратя на это n шагов. Тогда на первом уровне дерева будет находиться n узлов, каждый из них равен $\frac{n}{2}$. 

На следующем вызове рекурсий каждый из n узлов первого уровня вызовет $\frac{n}{2}$ подзадач размера $\frac{n}{4}$, затрачивая на это $\frac{n}{2}$ операций каждый. Тогда: общие затраты = количество узлов * шаги единичного разбиания = $n * \frac{n}{2} = \frac{n^2}{2}$ шагов. При этом на втором уровне дерева будет $n * \frac{n}{2} = \frac{n^2}{2}$ узлов, каждый размера $\frac{n}{4}$

Какая закономерность? С каждым последующим уровнем размер узла уменьшается вдвое. Количество узлов при этом увеличивается в такое количество раз, равное размеру узлов на предыдущем уровне. Количество шагов на уровне равно количеству узлов умноженное на размер единичного узла. Тогда:

На $k$ уровне: каждый узел имеет размер $\frac{n}{2^k}$ (пусть у нас корень - это 0 уровень). Тогда количество этих узлов составляет $n * \frac{n}{2} * \frac{n}{4} * ... * \frac{n}{2^{k-1}} = \frac{n^k}{2^{0+1+2+3...+k-1}} = \frac{n^k}{2^{\frac{k(k-1)}{2}}}$. И каждый узел будет выполнять разбиение за $\Theta(\frac{n}{2^k})$, то есть всего время разбиение уровня составит $n * \frac{n}{2} * \frac{n}{4} * ... * \frac{n}{2^{k-1}} * \Theta(\frac{n}{2^k}) = \frac{n^k}{2^{\frac{k(k-1)}{2}}} \Theta(\frac{n}{2^k}) = \Theta(\frac{n^{k+1}}{2^{\frac{k(k+1)}{2}}})$

Теперь вспомним, что наше дерево полное. От корня $n$ мы дойдем до узлов $1$, каждый раз уменьшая узел вдвое, за $\log_2 n$. 

Тогда время работы алгоритма $T(n)$ будет равно сумме делений каждого уровня. Найдем сумму: $\Theta(T(n)) = \sum_{k=0}^{\log_2 n} \Theta(\frac{n^{k+1}}{2^{\frac{k(k+1)}{2}}})$. Эта сумма будет определяться последним членом: $\frac{n^{log_2 n + 1}}{2^{\frac{\log_2n (\log_2n + 1)}{2}}} = n^{\frac{\log 2n}{\log 4}} = \Theta(n^{\log n})$

Ответ: $T(n) = \Theta(n^{\log n})$

\medskip

***************************************************************

\medskip

такое же достигается по мастер-теореме: $a = n, b =2, c = \log_2 n; n = O(n^{c - \varepsilon}) = O(n^{log_2 n - \varepsilon})$
верно $\forall \varepsilon \Rightarrow T(n) = \Theta (n^{c}) = \Theta (n^{\log n})$. Оно кажется правильным, но оценка $n$ через $O(n^{\log n})$... не знаю. странно

\bigskip


\end{document}